//323111047

\section{Introducción}

Planteamiento del problema. El reto técnico de esta actividad consiste en diseñar una aplicación de línea de comandos en Python 3 para resolver un laberinto representado mediante una matriz bidimensional. El problema central radica en la implementación de una función recursiva que sea capaz de llamarse a sí misma para explorar las direcciones de arriba, abajo, izquierda y derecha. Para lograr esto, es necesario identificar con precisión los casos base de éxito al encontrar la salida y de fallo al detectar una pared o un límite de la matriz, asegurando que cada llamada recursiva contribuya a encontrar una ruta válida desde la posición inicial marcada como S hasta la meta marcada como E.

\vspace{0.4cm}

Motivación. La necesidad de realizar este ejercicio surge de la importancia de comprender la recursividad como una técnica esencial para simplificar la resolución de problemas que presentan una naturaleza repetitiva o jerárquica. El pensamiento recursivo permite a un ingeniero diseñar soluciones compactas y lógicas para tareas que, de otro modo, requerirían estructuras iterativas sumamente complejas o difíciles de mantener. Estudiar la aplicación de funciones que se llaman a sí mismas dentro de un entorno de celdas permite visualizar el comportamiento de la pila de ejecución en tiempo real y entender la gestión de estados en la memoria, habilidades que son fundamentales para el estudio posterior de estructuras de datos avanzadas y la optimización de procesos computacionales.

\vspace{0.4cm}

Objetivos. El propósito fundamental de esta práctica es lograr la implementación funcional de la recursividad para la exploración sistemática de un laberinto en un entorno de Jupyter Notebook. Se busca que el alumno demuestre dominio en la construcción de funciones recursivas, validando el correcto funcionamiento de los casos base y la técnica de retroceso para limpiar caminos que no conducen a la salida. Finalmente, se pretende analizar la complejidad en tiempo y espacio de este enfoque recursivo, proporcionando una base técnica sólida para discutir la eficiencia del algoritmo y el cumplimiento de los requerimientos establecidos en la conclusión del reporte.